\section{Introduction}
\label{sec:intro}

Event cameras are bio-inspired sensors that report asynchronous per-pixel log-intensity changes, offering microsecond temporal resolution, high dynamic range, and no motion blur~\cite{gallego2020event}. These properties make them attractive for capturing fast dynamic 3D scenes. Meanwhile, 3D Gaussian Splatting (3DGS)~\cite{kerbl3Dgaussians} has emerged as a high-fidelity, real-time-renderable scene representation. A natural question arises: \emph{can we use a static event camera to live-update a known 3D Gaussian scene at high temporal resolution, reconstructing dynamic deformations online?}

This question is motivated by the \emph{digital-twin / scene-live-update} setting: a sensor is fixed in a known pose observing a 3D asset whose static, high-fidelity model $\Gzero$ has already been built (\eg\ by an offline RGB reconstruction). The asset then deforms or moves, and we wish to update the 3D model with the deformation at the highest possible temporal resolution, without re-capturing or re-reconstructing the full scene. A fixed event camera is the natural sensor for this task: it reports exactly the per-pixel log-intensity changes caused by the dynamic increment, at microsecond latency, with negligible bandwidth when the scene is static.

Existing event-based 3D works do not address this setting. EventNeRF~\cite{eventnerf} pioneered single-color-event NeRF reconstruction, while a rapidly growing line of work---Event-3DGS~\cite{event3dgs}, Event3DGS~\cite{xiong2025event3dgs}, EventSplat~\cite{eventsplat}, EvGGS~\cite{evggs}, E2GS~\cite{e2gs}, and EF-3DGS~\cite{ef3dgs}---replaces the implicit NeRF with explicit 3D Gaussians for faster, higher-fidelity rendering. Pose-free extensions such as IncEventGS~\cite{inceventgs} and EvTrajGS~\cite{evtrajgs} remove the need for known camera poses via SLAM-style tracking. Dynamic-scene variants---Dynamic EventNeRF~\cite{dynamiceventnerf}, Deformable NeRF with RGB and events~\cite{defnerfevent}, Event-Boosted Deformable 3DGS~\cite{eventboostdeform3dgs}, DiET-GS~\cite{dietgs}, EdMCGS~\cite{edmcgs}, FastEventDGS~\cite{fasteventdgs}, and ERF-GS~\cite{erfgs}---combine events with sparse RGB frames to model non-rigid motion. However, \emph{all of these methods reconstruct the entire scene from events (and possibly sparse frames), which limits fidelity because events carry only relative log-intensity changes, not absolute appearance}, and \emph{all assume a moving camera to generate parallax}.

We pursue the complementary and, to our knowledge, unexplored setting of a \emph{static} event camera observing a \emph{dynamic} scene, given a \emph{known high-fidelity baseline GS scene} $\Gzero$ captured beforehand. This setting has two key advantages that make high-temporal-resolution dynamic reconstruction tractable:
\begin{enumerate}[leftmargin=*,topsep=2pt,itemsep=1pt]
  \item \textbf{Static camera with known pose.} The pixel-to-Gaussian correspondence is fixed and precomputable: each pixel maps to a known set of Gaussians along its ray. Events, which are spatially sparse, thus localize directly to 3D Gaussians, enabling sparse, ray-localized 3D updates rather than dense image-wide reconstruction.
  \item \textbf{Known baseline scene.} Events need only explain the \emph{dynamic increment} $\dG(t)$ on top of $\Gzero$, not the entire scene. Appearance, geometry, and texture are provided by the baseline; events supply the high-frequency temporal deformation. This decouples fidelity (from $\Gzero$) from temporal resolution (from events).
\end{enumerate}

The physical bridge between events and Gaussians is the event-generation model: an event $e_k=(x_k,y_k,t_k,p_k)$ fires when
\begin{equation}
  \log I(x_k,y_k,t_k;\Gt) - \log I(x_k,y_k,t_{k-1};\Gt) = p_k\, C,
  \label{eq:event-model}
\end{equation}
where $I(\cdot;\Gt)$ is the rasterized intensity of the current Gaussian scene and $C$ is the contrast threshold. Equation~\ref{eq:event-model} is an observation equation relating 3D Gaussian parameters to 2D events through the differentiable rasterizer. This is the 3D analogue of event-to-video (E2VID)~\cite{rebecq2019e2vid}: where E2VID inverts the integral of Eq.~\ref{eq:event-model} to produce a 2D image per timestamp, we instead produce the \emph{3D deformation} that explains the observed events, keeping intensity on the scene side.

\medskip\noindent\textbf{The aperture problem.} A single event is one scalar constraint per pixel, but a Gaussian increment $\dmu_i\in\mathbb{R}^3$ has three degrees of freedom. The event loss therefore only constrains the component of $\dmu_i$ along the \emph{projected image gradient} direction; the tangential component (parallel to the edge) lies in the nullspace. This is the 3D analogue of the classic aperture problem~\cite{gallego2018contrast,barron1994aperture}, and we observe it empirically: the event-only loss admits a degenerate ``no-motion'' minimum. Prior event-3DGS works do not analyze this degeneracy because they rely on a moving camera and multi-view parallax to break it; in our static-camera setting it becomes the central challenge.

\medskip\noindent\textbf{Contributions.}
\begin{itemize}[leftmargin=*,topsep=2pt,itemsep=1pt]
  \item We formulate event-driven dynamic GS updating as a differentiable analysis-by-synthesis problem with an event-consistency loss that supervises both event and non-event pixels (Sec.~\ref{sec:method}), in a static-camera / known-baseline setting unexplored by prior event-3DGS work.
  \item We make the aperture degeneracy \emph{computable}: a per-Gaussian observability score from the structure tensor of the baseline render that (i) weights the event loss away from the degenerate minimum and (ii) triggers keyframes adaptively (Sec.~\ref{sec:obs}).
  \item We introduce a Contrast-Maximization regularizer on the predicted 2D Gaussian flow that provides a keyframe-free, edge-normal aperture-breaking signal in every window (Sec.~\ref{sec:cmax}).
  \item We instantiate the prototype with a differentiable orthographic rasterizer and a polynomial-motion event decoder, and evaluate on three synthetic benchmarks isolating the aperture problem (Sec.~\ref{sec:experiments}). On a pure straight-bar translation, event-only recovery is at chance ($0.53$ sign-agreement); the full model recovers correct-sign motion ($0.99$ with keyframe, $0.64$ with CMax alone), while the adaptive scheduler requests $\sim$$100\times$ fewer keyframes than fixed cadence.
\end{itemize}
